\chapter{Plant Nomenclature}

Following are some explanations of how taxonomy has been handled in this
book and some suggestions as to how taxonomy can be made more effective. The
Linneaus binomial system of classifying plants was based on the morphology of the
sexual parts of plants. One of its triumphs was that it led to a classification of plants
that was in accord with their evolutionary development. This accord was not only at
the macroscopic level but also at the molecular level: For example the alkaloid
berberine is characteristic of Berberidaceae, and there are many other examples of
particular chemicals and chemistries that are characteristic of plant families and
genera.

However, the Linneaus system had two misconceptions: At the time of
Linnaeus academic institutions still clung to latin as an emblem of intellectualism. This
is now archaic. A name like Eranthis hyemalis is an internationally recognized name
for a plant and belongs to no specific language. It is not a foreign word and is
specifically not a latin word so that to italicize it is grammatically incorrect. Plant
names \textit{have not been} italicized throughout this book, and it is anticipated that this
practice will become more prevalent with time. Plant names are not italicized in our
local newspaper or in ref. 1, and names are not italicized in other fields of science.

The second misconception was that plant species are immutable. The work of
Darwin and Darwinian Evolution showed the error in this. The problem is then how
can the Linnean system be so successful when its very premise is wrong. Recently
Gould has disussed this problein in a publication entitled "What is a Species" (ref. 48).
The answer is that species differentiate over a relatively short period of geological
time and remain distinct over long periods of geologic time. This means that at any
instant only a relatively few plants are undergoing species differentiation. Examples of
such active species differentiation at this present time would be the some of the small
Mediterranean orchids and the Phlox of Southeastern U.S. that center around Phlox
pulchra and Phlox glaberrima.

Plant families and genera have separated so far in an evolutionary sense that
they can be regarded as immutable. Thus a lily and a rose will never interbreed nor
exchange genetic information, and in effect Liliaceae and Rosaceae are for all
practical purposes permanent and immutable.

Species can then be defined by the principle that members of a species can
breed with others in the same species but not with individuals belonging to a different
species. The most obvious reason for lack of interbreeding is incompatibility
(intersterility) and the failure of crosses to set viable progeny. It also can be the result
of different flowering times and other differences in habits; Gould excludes
geographical separation as a failure to interbreed. These definitions are adopted in
this book. Genus can then be defined as a group that are interfertile but do not
normally interbreed. This definition is laborious to test, so that genera have been
delineated largely by inference. Yet it is remarkable that so few examples of
intergeneric hybrids have ever been made,and these few exceptions could be
eliminated by coalescing the genera involved. I can think of only one example at the
moment, and that is the hybrid of Belamcanda chinensis and Iris dichotoma.

All of the above seems clear, yet there are serious problems in taxonomy. The
most important function of taxonomy is to faciltate communication. This is more
important than taxonomy as a field in itself. Several practices have impeded this
objective. The principle that priority is given to the first given name seems simple at
first, but in fact has led to mischief. As one goes back in time the descriptions of plants
become more vague and ambiguous. At the present time and the present stage of the
situation, efforts to find an earlier description with a concomitant name change is
science at its worst, and such efforts should be ended right now. To resurrect some
earlier name at this stage hurts communication badly and serves no purpose other
than to get a \textit{cheap and trivial publication}. In chemistry an international conference
was convened in Geneva in 1892, and an end was put to further frivolous name
changes. Botanical taxonomy would be well advised to follow a similar route.

Names of plants should be changed only with the greatest reluctance and when
overpowering arguments favor such a change. This has recently come to the fore in
regard to the names of plant families. Nearly all the plant families have names derived
from a particular genus within the family, and the name ends in \textit{ceae}. However, there
are seven familes in which the name has been derived from a characteristic of the
family and the name ends in Ag. New names have been proposed for these seven
families which conform with more general practice. These new names and the old
name in parethesis are Apiaceae (Umbelliferae), Asteraceae (Compositae),
Brassicaceae (Cruciferae), Fabaceae (Leguminosae), Hypericaceae (Clusiasceae,
Guttiferae), Lamiaceae (Labiatae), and Poaceae (Graminae). These new names have
been adopted in this book primarily because it makes the ending of the name
diagnostic for a plant family, and thus further reduces any need for italics:
In virtually all fields of science the principle of Occam's razor is given first
priority. Applied to plant taxonomy this would state that no more plant names should
be invoked than are necessary for unambiguous and efficient communication. In the
past the practice of taxonomists in commemorating each other in plant names and the
desire for publications has often led to a needless proliferation of names. Fortunately
recent trends towards coalescing plant names is a step in the right direction. A name
like Paris quadrifolia conveys much more information than a commemorative name
like Tricyrtis baked. It is unfortunate that commemorative names ever got started. It is
obvious that commemorating influential people by naming a plant after them has been
used for financial and political benefit and to further ones career. All of this introduces
subjectivity and political influence into what should be an objective science.

